% Generated draft: requires mathematical and source review.
\documentclass[11pt]{article}
\usepackage[margin=25mm]{geometry}
\usepackage{fontspec}
\setmainfont{Latin Modern Roman}
\usepackage{amsmath,amssymb,mathtools,unicode-math}
\setmathfont{Latin Modern Math}
\usepackage{enumitem,needspace,xurl,hyperref}
\setlength{\parindent}{0pt}
\setlength{\parskip}{5pt}
\setlength{\emergencystretch}{3em}
\newcommand{\N}{\mathbb{N}}
\newcommand{\Z}{\mathbb{Z}}
\newcommand{\Q}{\mathbb{Q}}
\newcommand{\R}{\mathbb{R}}
\newcommand{\C}{\mathbb{C}}
\newcommand{\F}{\mathbb{F}}
\newcommand{\A}{\mathbb{A}}
\newcommand{\PP}{\mathbb{P}}
\newcommand{\E}{\mathbb{E}}
\newcommand{\eps}{\varepsilon}
\newcommand{\dd}{\,\mathrm{d}}
\begin{document}
\section*{UnsolvedMath ID: 800633}
\textbf{Source problem code:} catalogue entry
\subsection{Definitions and mathematical statement}
Let \(X\) be a topological space, and give every subset \(Y\subseteq X\) the subspace topology. Under the usual cover conventions, a Borel \(\omega\)-cover of \(X\) is a countable family \(\mathcal U\) of Borel subsets of \(X\) such that
\[
X\notin\mathcal U
\quad\text{and}\quad
\forall F\subseteq X\;\bigl(F\text{ finite}\Longrightarrow\exists U\in\mathcal U\;(F\subseteq U)\bigr).
\]
Write \(\mathcal B_\Omega(X)\) for the collection of these covers. Borel sets and covers of a subspace \(Y\) are understood relative to \(Y\).

The property \(\mathsf S_1(\mathcal B_\Omega,\mathcal B_\Omega)\) holds for \(X\) if, for every sequence \((\mathcal U_n)_{n\in\N}\) from \(\mathcal B_\Omega(X)\), there are choices \(U_n\in\mathcal U_n\) such that
\[
\{U_n:n\in\N\}\in\mathcal B_\Omega(X).
\]
The property \(\mathsf S_{\mathrm{fin}}(\mathcal B_\Omega,\mathcal B_\Omega)\) holds for \(X\) if, for every such sequence, there are finite subfamilies \(\mathcal F_n\subseteq\mathcal U_n\) such that
\[
\bigcup_{n\in\N}\mathcal F_n\in\mathcal B_\Omega(X).
\]
Here the union is a family of selected Borel sets.

For a property \(P\), hereditariness means that every space under consideration with \(P\) has \(P\) on each of its subspaces. The two questions are whether one can prove, separately for
\[
P\in\left\{\mathsf S_1(\mathcal B_\Omega,\mathcal B_\Omega),\ \mathsf S_{\mathrm{fin}}(\mathcal B_\Omega,\mathcal B_\Omega)\right\},
\]
the implication
\[
\forall X\;\left(P(X)\Longrightarrow\forall Y\subseteq X\;P(Y)\right),
\]
where \(X\) ranges over the source's intended class of spaces. The supplied excerpt does not specify that class or the foundational theory intended by ``provably.'' Thus the displayed formulation retains an unresolved domain qualification; it does not silently replace the intended class by all topological spaces. The cover definitions above use standard conventions that also require checking against the full source.

For context, the supplied source reports that none of the properties involving open covers considered in that paper is hereditary, whereas \(\mathsf S_1(\mathcal B,\mathcal B)\) and all properties of the form \(\Pi(\mathcal B_\Gamma,\mathcal B)\) are hereditary, citing [ideals]. It also reports that \(\mathsf S_1(\mathcal B_\Omega,\mathcal B_\Gamma)\) is not hereditary, citing [MilNonGamma]. Here \(\mathcal B\) and \(\mathcal B_\Gamma\) conventionally denote countable Borel covers and countably infinite Borel \(\gamma\)-covers, respectively; in a \(\gamma\)-cover each point belongs to all but finitely many members. The excerpt does not specify the range of the schema \(\Pi\). These contextual assertions are supplied source claims, not independently verified conclusions.
\subsection{Short English statement}
For each of the two Borel \(\omega\)-cover selection properties, ask whether every subspace of a space having that property must also have it. The first property selects one set from each cover; the second selects finitely many. In either case, the selected sets must together form an \(\omega\)-cover. The intended class of spaces and the meaning of provability remain unspecified in the supplied excerpt. The registry labels the problem open; this is a supplied status, not an independently verified current assessment.
\subsection{Sources}
\begin{itemize}
\item Local UnsolvedMath registry, supplied complete record: permanent ID \(800633\), exact source problem code \texttt{catalogue entry}, AMR Open Problem Lists. The record attributes the problem to Tsaban, with proposed year 2003, and supplies status ``open.'' Its creation and update timestamps are both \texttt{2026-10-16T14:00:00Z}.
\item Tsaban, \emph{Selection Principles in Mathematics: A Milestone of Open Problems} (2003), Problem 5. Supplied source URL: \url{https://arxiv.org/abs/math/0312182}. The record reports access on 2026-10-16 and extraction by \texttt{arxiv\_latex}; these are supplied provenance details, not a claim of independent source access or verification.
\item Supplied attribution and context: [ideals, Problem 4] and [MilNonGamma, Question 6]. The contextual hereditary and nonhereditary assertions are attributed in the excerpt to [ideals] and [MilNonGamma], respectively. Full bibliographic entries for these labels were not supplied.
\item Supplied upstream identity: \(800633\); upstream revision \texttt{d371ee917ac9647fbb06cd6fc529deec6f3dda82}.
\end{itemize}
\end{document}
